\section{未来路线：Generalist Agent}

\subsection{三进三出的全模态模型}

Qwen 的终极目标是一个能够同时\textbf{理解和生成}文本、视觉、音频三种模态的统一模型：

\begin{itemize}
    \item 已实现：文本理解/生成、视觉理解、音频理解/生成
    \item 待实现：将视觉生成融入统一模型
    \item 终极形态：``三进三出''——三种模态自由输入输出
\end{itemize}

\subsection{Multi-turn RL with Environment Feedback}

\begin{importantbox}{Long Horizon Reasoning}
林俊旸认为 RL 的未来不是做 DeepSeek R1 式的数学竞赛自对话，而是 \textbf{Multi-turn RL with Environment Feedback for Long Horizon Reasoning}：
\begin{itemize}
    \item 模型在真实环境中执行长链任务（如花两天完成人类两个月的工作）
    \item 通过环境反馈进行学习，而非仅靠自我对话
    \item 适用于虚拟世界（Digital Agent）和物理世界（Embodied Agent）
    \item 本质上是将 Natural Language Instruction 转换为 Executable Action
\end{itemize}
\end{importantbox}

\begin{knowledgebox}{RL Compute 的趋势}
XAI 的宣传显示其 RL 计算量已接近 Pre-training 计算量。虽然林俊旸认为``有点浪费''，但这确实说明 RL 有巨大的想象空间。Qwen 团队更关心的不是让模型成为最强数学大脑，而是让它能像真实的人一样为生物和社会做贡献。
\end{knowledgebox}

\subsection{从 Digital Agent 到 Physical Agent}

\begin{itemize}
    \item \textbf{Digital Agent}：同时操作 GUI 和 API，是当前最成熟的方向
    \item \textbf{Physical Agent}：操控话筒、倒水等物理世界任务，需要 VLA（Vision-Language-Action）支持
    \item 统一视角：无论 Coding Agent 还是 VLA，本质都是``将自然语言指令转化为可执行动作''
\end{itemize}

\subsection{本章小结}

Qwen 的技术路线清晰：全模态统一 → Agent RL → Digital/Physical Agent。这是一条从``理解世界''到``改变世界''的进化之路。

